%%%%%%%%%%%%%%%%%%%%%%%%%%%%%%%%%%%%%%%%%%%%%%%%%%%%%%%%%%%%%%%%%%%%%%%%%%%%%%%%
% ICRA 2027 submission skeleton.
% Rules (https://2027.ieee-icra.org/contribute/call-for-icra-2027-papers-now-accepting-submissions/):
%   - 8 pages max INCLUDING references; over-length papers are rejected without review
%   - double-anonymous: no author names/affiliations/identifying acknowledgments
%   - ieeeconf double-column, US letter, PDF only; optional video <= 180 s, <= 20 MB
%   - AI-assisted writing must be disclosed in the acknowledgment section
%   - submission deadline: September 15, 2026 (23:59 PST)
% Every \todo{} is a hole for the robotics co-authors or for results to be filled.
%%%%%%%%%%%%%%%%%%%%%%%%%%%%%%%%%%%%%%%%%%%%%%%%%%%%%%%%%%%%%%%%%%%%%%%%%%%%%%%%

\documentclass[letterpaper, 10 pt, conference]{ieeeconf}

\IEEEoverridecommandlockouts
\overrideIEEEmargins

\usepackage{amsmath,amssymb}
\usepackage{graphicx}
\usepackage{booktabs}
\usepackage{xcolor}
\usepackage[hidelinks]{hyperref}

% drafting helpers: \todo{} flags a hole, \draft{} flags prose that needs a rewrite
\newcommand{\todo}[1]{\textcolor{red}{[TODO: #1]}}
\newcommand{\draft}[1]{\textcolor{blue}{#1}}

% notation
\newcommand{\R}{\mathbb{R}}
\newcommand{\Hs}{\mathcal{H}}
\newcommand{\Fs}{\mathcal{F}}
\newcommand{\EI}{\mathrm{EI}}
\newcommand{\vq}{\mathbf{q}}
\newcommand{\vf}{\mathbf{f}}
\newcommand{\vtau}{\boldsymbol{\tau}}

\title{\LARGE \bf
Learning to Manipulate a Whip with Functional Bayesian Optimization
}

% Double-anonymous: keep the author block anonymous until camera-ready.
\author{Anonymous Submission}

\begin{document}

\maketitle
\thispagestyle{empty}
\pagestyle{empty}

%%%%%%%%%%%%%%%%%%%%%%%%%%%%%%%%%%%%%%%%%%%%%%%%%%%%%%%%%%%%%%%%%%%%%%%%%%%%%%%%
\begin{abstract}
Striking a target with a whip is a canonical example of dynamic manipulation of
a deformable object: a 4-degree-of-freedom (DOF) arm drives a 50-DOF flexible
chain through a brief, precisely timed motion. Nah~\emph{et al.}~\cite{nah2023whip}
showed in simulation that this task can be solved by restricting the arm to a
single minimum-jerk submovement and searching its nine parameters with a
derivative-free optimizer, hitting five of six targets within 39 to 249
simulated episodes. We revisit the same task and ask whether the restriction to
minimum-jerk primitives is necessary. We treat the zero-torque joint trajectory
itself as the decision variable and optimize it with functional Bayesian
optimization (FBO): the trajectory is represented as a sparse mixture of
radial basis functions, a Gaussian process over trajectories is fitted through
their distance in a reproducing kernel Hilbert space (RKHS), and expected
improvement is maximized jointly over basis locations, widths and amplitudes.
\todo{results sentence: number of episodes to hit each target, whether target 6
is hit, comparison to the 39--249 episodes of DIRECT-L}. The learned
trajectories are not minimum-jerk, \todo{one sentence on what they look like},
which suggests that sample-efficient search over trajectory space, rather than
a fixed motor primitive, is what makes the problem tractable.
\end{abstract}

%%%%%%%%%%%%%%%%%%%%%%%%%%%%%%%%%%%%%%%%%%%%%%%%%%%%%%%%%%%%%%%%%%%%%%%%%%%%%%%%
\section{INTRODUCTION}

\draft{
Manipulating a whip confronts a controller with everything that makes
deformable objects hard: a high-dimensional state, nonlinear wave propagation
along the object, and extreme sensitivity of the outcome to the timing of the
input~\cite{krotov2022whip,nah2023whip}. Yet humans learn to crack whips and
strike targets with them, which has motivated the hypothesis that human motor
control side-steps the complexity by composing a small number of dynamic
primitives~\cite{hogan2012primitives,flash1985minjerk}.
}

\draft{
Nah~\emph{et al.}~\cite{nah2023whip} tested this hypothesis in simulation. A
4-DOF upper limb model, driven by a joint-space impedance controller, holds a
25-node lumped-parameter whip. The arm follows a single rest-to-rest
minimum-jerk motion in joint space, so the whole action is described by nine
numbers: the initial posture, the final posture and the movement duration. The
DIRECT-L global optimizer~\cite{gablonsky2001direct} found parameters that hit
five of six targets in fewer than 250 episodes, and came within 10~cm of the
sixth. The authors attribute this efficiency to the primitive: the search space
is nine-dimensional instead of a trajectory space with thousands of samples.
}

\draft{
This leaves two questions open. First, is the minimum-jerk constraint necessary
for tractability, or merely sufficient? A richer trajectory class could reach
targets that a single submovement cannot, provided the search remains
sample-efficient. Second, is DIRECT-L the right tool? Each episode is an
expensive simulation (or, on hardware, a physical trial), which is the regime
Bayesian optimization (BO) is designed for~\cite{shahriari2016bo,frazier2018bo}.
}

\draft{
We answer both questions with functional Bayesian optimization: BO whose
decision variable is a function rather than a vector. The zero-torque joint
trajectory $\vq_0(t)$ is represented as a sparse mixture of radial basis
functions with free centers, widths and amplitudes. A Gaussian process over
trajectories is fitted using the RKHS distance between mixtures, which has a
closed form, and expected improvement is maximized by gradient ascent over all
mixture parameters. The basis size grows during the run so early episodes are
spent on coarse trajectories and later ones on refinements.
}

\noindent\textbf{Contributions.}
\begin{itemize}
  \item We reproduce the whip-targeting benchmark of~\cite{nah2023whip}
  (4-DOF arm, 50-DOF whip, six targets, joint-space impedance control) as a
  differentiable simulator and release it \todo{link at camera-ready}.
  \item We show that functional BO hits \todo{$n$ of six} targets in
  \todo{$k$} episodes, \todo{$x\times$} fewer than DIRECT-L on the same task.
  \item We show that the minimum-jerk restriction is not needed: the trajectories
  found by FBO are \todo{characterize}, and \todo{target 6 result}.
\end{itemize}

%%%%%%%%%%%%%%%%%%%%%%%%%%%%%%%%%%%%%%%%%%%%%%%%%%%%%%%%%%%%%%%%%%%%%%%%%%%%%%%%
\section{RELATED WORK}

\subsection{Whip manipulation and dynamic primitives}
\draft{
Krotov~\emph{et al.}~\cite{krotov2022whip} recorded humans hitting a target
with a bullwhip and found a bell-shaped hand speed profile in both discrete and
rhythmic styles. Nah~\emph{et al.} built on this in a series of simulation
studies~\cite{nah2020whip,nah2021whip,nah2023whip}, culminating in the
targeting task we adopt. Their controller is a joint-space impedance controller
\cite{hogan1985impedance} following a minimum-jerk virtual
trajectory~\cite{flash1985minjerk}; the primitive is fixed and only its
parameters are optimized.
}
\todo{robotics co-authors: other deformable/dynamic manipulation work worth citing
(cables, cloth, rope, ``robots of the lost arc'', iterative residual policy)}

\subsection{Policy search with few trials}
\draft{
Parametrized motion primitives such as dynamic movement primitives
\cite{ijspeert2013dmp} are the standard way to make policy search
sample-efficient in robotics~\cite{kober2011policysearch,chatzilygeroudis2020survey}.
Bayesian optimization has been used to tune gait parameters
\cite{calandra2016gaits} and controller gains directly on hardware in tens of
trials. In all these cases the policy class is fixed in advance and BO searches a
low-dimensional parameter vector; we instead search the trajectory itself.
}

\subsection{Bayesian optimization over functions}
\draft{
Functional BO extends the surrogate from vectors to elements of a function
space. Vien~\emph{et al.}~\cite{vien2018fbo} optimize over an RKHS with a
kernel on function norms; Vellanki~\emph{et al.}~\cite{vellanki2019fbo} add
shape priors; Shilton~\emph{et al.}~\cite{shilton2020fbo} discretize the
domain on a grid. Our method \cite{ours2026fbo} searches a sparse sub-manifold
of the RKHS by jointly optimizing basis point locations and coefficients, which
keeps the acquisition problem low-dimensional while letting the basis adapt to
the objective. This paper is the first application of that method to a
physical control task with a strong published baseline.
}

%%%%%%%%%%%%%%%%%%%%%%%%%%%%%%%%%%%%%%%%%%%%%%%%%%%%%%%%%%%%%%%%%%%%%%%%%%%%%%%%
\section{THE WHIP TARGETING TASK}
\label{sec:task}

We follow the model of~\cite{nah2023whip}; parameters are listed in
Table~\ref{tab:params}.

\subsection{Arm and whip model}
\draft{
The upper limb is a two-link open chain with a 3-DOF shoulder (flexion/extension
$J_1$, adduction/abduction $J_2$, lateral/medial rotation $J_3$) and a 1-DOF
elbow ($J_4$). Segment lengths, masses and inertias follow Hatze's
anthropometric model. The whip is a lumped-parameter chain of $N=25$ point
masses joined by massless rods and 2-DOF universal joints, each joint carrying
a linear rotational spring $k$ and damper $b$; the joint attaching the whip to
the hand has $k=b=0$. The complete system has $4 + 2N = 54$ DOF.
}
\todo{figure: schematic of arm, whip, joint axes, and the six targets (cf.\ Fig.~1 and Fig.~11 of~\cite{nah2023whip})}

\subsection{Impedance controller}
The arm is driven by a first-order joint-space impedance controller with gravity
compensation,
\begin{equation}
  \vtau = K_q\,(\vq_0 - \vq) + B_q\,(\dot\vq_0 - \dot\vq) + \vtau_G,
  \label{eq:impedance}
\end{equation}
where $\vq\in\R^4$ is the measured joint configuration, $\vq_0(t)$ is the
zero-torque (virtual) trajectory, $K_q$ is the measured human joint stiffness of
\cite{nah2023whip} and $B_q = 0.05\,K_q$. The gravity torque $\vtau_G$
compensates the arm links and the whip mass lumped at the hand.

\subsection{Objective}
\draft{
An episode simulates $T=2.5$~s from rest with the whip hanging. The cost is the
closest approach of the whip tip to the target, minus a contact radius $r$,
\begin{equation}
  L(\vq_0) = \max\Big(0,\; \min_{t,\,i\in\mathcal{I}_{\text{tip}}}
  \big\|\mathbf{p}_i(t) - \mathbf{p}^{\star}\big\| - r\Big),
  \label{eq:objective}
\end{equation}
where $\mathbf{p}_i(t)$ is the position of node $i$ and
$\mathcal{I}_{\text{tip}}$ indexes the last five nodes. A hit is $L=0$.
}
\todo{align with the paper: they use all 25 nodes and a 3~cm target radius; we
use the last 5 nodes and $r=5$~cm. Decide and state which, and justify (tip-only
prevents ``cheating'' hits with the handle end).}

\subsection{Targets}
Six targets are placed on spheres around the shoulder: targets 1--3 at radius
$R - d$ and targets 4--6 at $R$, where $R = L_1 + L_2 + N l = 2.385$~m is the
combined reach of arm and whip and $d=0.4$~m. Azimuth/elevation are
$(0^\circ,0^\circ)$, $(45^\circ,0^\circ)$ and $(45^\circ,45^\circ)$ in each
group. Targets 4--6 can be reached only by the tip at full extension.

\subsection{Baseline: minimum-jerk primitive with DIRECT-L}
\draft{
In~\cite{nah2023whip} the virtual trajectory is a single rest-to-rest
minimum-jerk motion
\begin{equation}
  \vq_0(t) = \vq_{0,i} + (\vq_{0,f} - \vq_{0,i})\,
  \big(10\,s^3 - 15\,s^4 + 6\,s^5\big),\quad s = \min(t/D, 1),
  \label{eq:minjerk}
\end{equation}
with nine parameters $(\vq_{0,i}, \vq_{0,f}, D)$ optimized by DIRECT-L
\cite{gablonsky2001direct} for at most 600 iterations, halting on the first
hit. The reported episode counts to the first hit are 39--249 for targets 1--5;
target 6 was not hit (closest approach 10~cm).
}

\begin{table}[t]
\caption{Model parameters, from Table~3 of~\cite{nah2023whip}.}
\label{tab:params}
\centering
\begin{tabular}{llr}
\toprule
Quantity & Symbol & Value \\
\midrule
Upper arm / forearm length & $L_1, L_2$ & 0.294, 0.291 m \\
Upper arm / forearm mass & $M_1, M_2$ & 1.595, 0.869 kg \\
Whip nodes & $N$ & 25 \\
Node mass & $m$ & 0.012 kg \\
Segment length & $l$ & 0.072 m \\
Joint stiffness / damping & $k, b$ & 0.242 N$\cdot$m/rad, 0.092 N$\cdot$m$\cdot$s/rad \\
Contact radius & $r$ & \todo{0.05} m \\
Horizon / time step & $T, \Delta t$ & 2.5 s, 0.5 ms \\
\bottomrule
\end{tabular}
\end{table}

%%%%%%%%%%%%%%%%%%%%%%%%%%%%%%%%%%%%%%%%%%%%%%%%%%%%%%%%%%%%%%%%%%%%%%%%%%%%%%%%
\section{FUNCTIONAL BAYESIAN OPTIMIZATION}
\label{sec:method}

\todo{compress from the FBO paper~\cite{ours2026fbo}; keep only what a robotics
reader needs to trust the results}

\subsection{Trajectories as sparse RKHS elements}
\draft{
We parametrize the zero-torque trajectory through a vector-valued function
$\vf:[0,1]\to\R^4$ of normalized time $s = t/T$,
\begin{equation}
  \vq_0(t) = \mathbf{c} + \mathbf{A} \odot \vf(t/T),
  \label{eq:reference}
\end{equation}
where $\mathbf{c}$ and $\mathbf{A}$ center and scale each joint to its range
of motion, so $\vf$ is dimensionless. Each output of $\vf$ is a mixture of $m$
Gaussian radial basis functions,
\begin{equation}
  f_j(s) = \sum_{i=1}^{m} a_{ij}\,
  \exp\!\Big(-\frac{(s - x_{ij})^2}{2\,\ell_{ij}^2}\Big),
  \label{eq:mixture}
\end{equation}
with centers $x_{ij}\in[0,1]$, widths $\ell_{ij}$ and amplitudes
$a_{ij}\in[-1,1]$ all free. A trajectory is therefore a point $\theta\in\R^{3mk}$
with $k=4$, and the set of such mixtures is a sub-manifold of the RKHS $\Hs$ of
a wider Gaussian kernel. The minimum-jerk primitive of \eqref{eq:minjerk} is not
in this family, and conversely a single atom ($m=1$) already allows a
non-monotone, return-to-start motion.
}

\subsection{Gaussian process over trajectories}
\draft{
Rather than a GP on $\theta$, which would be invariant to nothing (permuting
atoms changes $\theta$ but not the trajectory), we place a GP on $\vf$ itself
with a stationary kernel in RKHS distance,
\begin{equation}
  \kappa(\vf, \vf') = \nu\,\phi\Big(\sqrt{\textstyle\sum_j \rho_j\,
  \|f_j - f'_j\|_{\Hs_{\ell_0}}^2}\Big),
  \label{eq:kernel}
\end{equation}
where $\phi$ is a Mat\'ern-5/2 profile, $\rho_j$ weights the joints, and
$\|\cdot\|_{\Hs_{\ell_0}}$ is the norm of the Gaussian RKHS with ambient
lengthscale $\ell_0$. For RBF mixtures this norm has a closed form, since the
inner product of two Gaussians in a Gaussian RKHS is again a Gaussian in their
centers. The hyperparameters $(\ell_0, \rho, \nu)$ and a constant trend are
fitted by marginal likelihood, warm-started between iterations.
}

\subsection{Acquisition and basis growth}
\draft{
We maximize (log) expected improvement~\cite{ament2023logei} by multi-start
gradient ascent over $\theta$: candidates are drawn by Latin hypercube over
$(\log\ell, x, a)$ and around the best observed trajectories, and the best
$n$ starts are refined with L-BFGS. The basis size follows the schedule
$m = 1, 2, 4, \dots$: at each stage boundary every atom is split in two with
half amplitude, which leaves all observed trajectories unchanged while doubling
the search dimension. Early episodes are thus spent on single-pulse motions and
later ones on refinements.
}
\todo{algorithm box: initial LHS design, loop (fit GP, maximize EI, simulate, record), stage boundaries}

%%%%%%%%%%%%%%%%%%%%%%%%%%%%%%%%%%%%%%%%%%%%%%%%%%%%%%%%%%%%%%%%%%%%%%%%%%%%%%%%
\section{EXPERIMENTS}
\label{sec:experiments}

\subsection{Protocol}
\draft{
For each of the six targets we run FBO with a budget of 128 episodes, of which
the first $n_0$ form a Latin hypercube design in the $m=1$ family. We compare
three basis schedules: fixed $m=1$ ($n_0=64$), growing to $m=2$ ($n_0=32$) and
growing to $m=4$ ($n_0=16$); the initial design is sized so every stage
doubles the number of observations. \todo{seeds: currently seed 0 only; run
$\geq 5$ seeds and report median and interquartile range.} The simulator is a
semi-implicit Euler integrator with implicit joint damping at $\Delta t = 0.5$~ms,
matching MuJoCo's integrator~\cite{todorov2012mujoco}; \todo{report validation
of our simulator against the published Table~1 optima (closest approach when
replaying the paper's nine parameters).}
}

\subsection{Sample efficiency}
\todo{Fig.: best-so-far closest approach vs.\ episode for each target, FBO
($m=1,2,4$) against the DIRECT-L curves of Fig.~2 in~\cite{nah2023whip}; mark
the first-hit episode.}

\todo{Table: episodes to first hit per target, FBO vs.\ DIRECT-L (39, 78, 39,
$\sim$170, 249, miss). Include final closest approach for misses.}

\begin{table}[t]
\caption{Episodes to first hit. DIRECT-L values from~\cite{nah2023whip}.}
\label{tab:episodes}
\centering
\begin{tabular}{lcccc}
\toprule
Target & DIRECT-L & FBO $m{=}1$ & FBO $m{=}2$ & FBO $m{=}4$ \\
\midrule
1 & \todo{} & \todo{} & \todo{} & \todo{} \\
2 & \todo{} & \todo{} & \todo{} & \todo{} \\
3 & 39 & \todo{} & \todo{} & \todo{} \\
4 & \todo{} & \todo{} & \todo{} & \todo{} \\
5 & 249 & \todo{} & \todo{} & \todo{} \\
6 & miss (0.10 m) & \todo{} & \todo{} & \todo{} \\
\bottomrule
\end{tabular}
\end{table}

\subsection{Beyond minimum jerk}
\todo{Fig.: learned $\vq_0(t)$ per joint for each target, overlaid on the
paper's minimum-jerk optimum; hand path in task space. Quantify the departure
from minimum jerk (e.g.\ number of speed peaks, jerk cost ratio).}

\todo{Fig.: target 6 rollout if hit; discuss what the extra freedom buys.}

\subsection{Ablations}
\todo{choose 2--3: (i) basis schedule fixed vs.\ growing; (ii) GP on
$\theta$ (vector BO) vs.\ GP on $\vf$ (RKHS kernel), same parametrization;
(iii) DIRECT-L on the same RBF parametrization, to separate parametrization
from optimizer; (iv) BO on the nine min-jerk parameters, to separate optimizer
from primitive.}

\subsection{Sensitivity}
\todo{optional: replicate the duration-sensitivity analysis of~\cite{nah2023whip}
on the FBO solutions; is the FBO trajectory more or less robust to timing
perturbation?}

%%%%%%%%%%%%%%%%%%%%%%%%%%%%%%%%%%%%%%%%%%%%%%%%%%%%%%%%%%%%%%%%%%%%%%%%%%%%%%%%
\section{DISCUSSION}
\todo{robotics co-authors: implications for real-robot trials (episode cost,
safety of exploratory trajectories, impedance on hardware); limitations
(simulation only, single seed, arm model without wrist); what FBO does not
address (sim-to-real).}

%%%%%%%%%%%%%%%%%%%%%%%%%%%%%%%%%%%%%%%%%%%%%%%%%%%%%%%%%%%%%%%%%%%%%%%%%%%%%%%%
\section{CONCLUSION}
\draft{
We revisited the whip targeting benchmark of Nah~\emph{et al.} and replaced the
fixed minimum-jerk primitive and DIRECT-L search with functional Bayesian
optimization over the virtual trajectory. \todo{one sentence of results}. The
result supports a different reading of the original study: what makes the task
learnable in few trials is a search that exploits smoothness of the objective
over trajectory space, and a fixed primitive is one way but not the only way to
get it.
}

% ICRA: balance columns on the last page; adjust once content is final
%\addtolength{\textheight}{-12cm}

%%%%%%%%%%%%%%%%%%%%%%%%%%%%%%%%%%%%%%%%%%%%%%%%%%%%%%%%%%%%%%%%%%%%%%%%%%%%%%%%
\section*{ACKNOWLEDGMENT}
\todo{camera-ready only (anonymous at submission). ICRA 2027 requires
disclosure of any AI-assisted writing, naming the system and the sections.}

%%%%%%%%%%%%%%%%%%%%%%%%%%%%%%%%%%%%%%%%%%%%%%%%%%%%%%%%%%%%%%%%%%%%%%%%%%%%%%%%
\bibliographystyle{ieeetr}
\bibliography{references}

\end{document}
